% Generated draft: requires mathematical and source review.
\documentclass[11pt]{article}
\usepackage[margin=25mm]{geometry}
\usepackage{fontspec}
\setmainfont{Latin Modern Roman}
\usepackage{amsmath,amssymb,mathtools,unicode-math}
\setmathfont{Latin Modern Math}
\usepackage{enumitem,needspace,xurl,hyperref}
\setlength{\parindent}{0pt}
\setlength{\parskip}{5pt}
\setlength{\emergencystretch}{3em}
\newcommand{\N}{\mathbb{N}}
\newcommand{\Z}{\mathbb{Z}}
\newcommand{\Q}{\mathbb{Q}}
\newcommand{\R}{\mathbb{R}}
\newcommand{\C}{\mathbb{C}}
\newcommand{\F}{\mathbb{F}}
\newcommand{\A}{\mathbb{A}}
\newcommand{\PP}{\mathbb{P}}
\newcommand{\E}{\mathbb{E}}
\newcommand{\eps}{\varepsilon}
\newcommand{\dd}{\,\mathrm{d}}
\begin{document}
\section*{UnsolvedMath ID: 10400067}
\textbf{Source problem code:} catalogue entry
\subsection{Definitions and mathematical statement}
Problem 3.17 asks for a topological construction of the knot invariant denoted \(P_K^\theta\), the 2-loop polynomial. In the usual perturbative-diagram terminology, loop degree refers to the first Betti number of the connected graph off the knot skeleton. The theta graph has two vertices joined by three edges and first Betti number two. This is distinct from Vassiliev degree two; the 2-loop part can contain contributions in many diagram degrees.

The requested construction should extract \(P_K^\theta\) from specified topological data of the knot and identify the result with the source's polynomial. The excerpt supplies neither the defining 2-loop expression nor its variables, coefficient ring, denominator-clearing conventions or normalization. Consequently a particular numerator formula in terms of Alexander polynomials, an edge-variable convention, or a specialization cannot be chosen as its definition from this record.

``Topological construction'' is also not given a formal criterion. The question is retained as a construction programme for the full named polynomial, not reduced to computing a low-degree coefficient or proving that an existing expression is an isotopy invariant. No prohibition on analytic ingredients or requirement of a particular surgery, surface or clasper construction is added.
\subsection{Short English statement}
Construct the complete 2-loop knot polynomial \(P_K^\theta\) by a topological method. Its precise definition and normalization, and the intended construction criterion, are missing; 2-loop degree must not be confused with finite-type degree two.
\subsection{Sources}
\begin{itemize}
\item Local UnsolvedMath registry, permanent ID 10400067, source code catalogue entry; no proposer or proposal year supplied; updated 2026-08-06.
\item Ohtsuki, \emph{Problems on invariants of knots and 3-manifolds} (2002), Problem 3.17, PDF page 67: \url{https://msp.org/gtm/2002/04/gtm-2002-04-024s.pdf}. Registry access date 2026-07-29; PDF font-structure extraction. Provenance distinguishes 231 numbered bold headings. Present primary-literature status and redistribution terms require review.
\item The supplied summary repeats the topological-construction request without a defining formula, exact further reference, or verified resolution.
\end{itemize}
\end{document}
